\subsection{Drop Down Dew}\label{sec:fourth-week}

On the Fourth Sunday the station is the church of the Twelve Apostles, and the
Mass begins with the whole of Isaiah's verse, \latin{Rorate, cæli, desuper}:

\begin{quote}
Drop down dew, ye heavens, from above, and let the clouds rain the just: let
the earth be opened, and bud forth a saviour: and let justice spring up
together: I the Lord have created him.\par
{\smallskip\noindent\small ---Isaiah 45:8\par}
\end{quote}

The prophet's two images stand side by side: the Just One coming down from heaven like dew, and a Saviour springing up out of
the earth like a plant. He comes from above and from below, from God and from
us.

Scripture has another dew. Gideon, before he
went out to save Israel, asked God for a sign:

\begin{quote}
I will put this fleece of wool on the floor: if there be dew on the fleece
only, and it be dry on all the ground beside, I shall know that by my hand, as
thou hast said, thou wilt deliver Israel. And it was so. And rising before day,
wringing the fleece, he filled a vessel with the dew.\par
{\smallskip\noindent\small ---Judges 6:37--38\par}
\end{quote}

St Bernard, preaching on the Annunciation, saw the Virgin in that fleece:

\begin{quote}
Other references might be quoted equally suitable to the Virgin Mother and to
the Son of God---Gideon's fleece, for instance, cut from the flesh without
wounding it, and placed on the dry ground, where the dry fleeces are in turn
moistened by the dew. This similitude represents the flesh assumed from the
flesh of Mary without injury to her virginity. Upon it Heaven dropped down dew,
filling it with the plenitude of the Divinity, and from that fullness we have
all received---we who, but for it, were as parched and arid soil. The Psalmist
seems to refer very beautifully to this fact in Gideon's history. In Ps. lxxi.
we read: ``He shall descend like rain upon the fleece, and as showers falling
gently upon the earth.'' Gently at first, and without noise of human operation,
He fell softly into the Virgin's womb; afterwards, when the Apostles announced
Him, it was with the noise of words and the display of miracles.\par
{\smallskip\noindent\small ---St Bernard, Second Sermon on the Virgin Mother\par}
\end{quote}

St Cyril of Jerusalem had said the same of the first coming: it was ``like rain
on a fleece'', unobserved. And on 1 January, the eighth day of Christmas, the
Church will sing it to the Child himself, in an antiphon at Lauds whose Latin,
like the psalm Bernard quoted, has rain on the fleece, \latin{sicut pluvia in vellus}, and
whose English has dew: ``When, in an unspeakable manner, Thou wast born of a
Virgin, the Scriptures were fulfilled; to save man, Thou camest down as dew upon
the fleece: we praise Thee, O our God!''

\subsubsection*{The fifteenth year of Tiberius}

The Gospel of this Sunday, the same that was read on the Ember Saturday, begins
like a page of history:

\begin{quote}
Now in the fifteenth year of the reign of Tiberius Caesar, Pontius Pilate being
governor of Judea, and Herod being tetrarch of Galilee, and Philip his brother
tetrarch of Iturea and the country of Trachonitis, and Lysanias tetrarch of
Abilina: Under the high priests Annas and Caiphas: the word of the Lord was
made unto John, the son of Zachary, in the desert.\par
{\smallskip\noindent\small ---Luke 3:1--2\par}
\end{quote}

St Gregory the Great's homily on this Gospel is headed as preached on an
Ember Saturday, and it is read at Matins of the Ember Saturday still. He begins with the
rulers: ``The date, at which the Fore-runner of our Redeemer entered on
his public office of preaching, is indicated to us by the name of the ruler of
the Roman Commonwealth, and by those of the princes of Palestine.'' The dew
falls into a year that can be counted, under rulers who can be named, and one
of them, Pontius Pilate, is the governor to whom the Lord will one day be brought
bound.

Into that year John comes ``preaching the baptism of penance for the remission
of sins'', and he cries Isaiah's words: ``Prepare ye the way of the Lord, make
straight his paths. Every valley shall be filled and every mountain and hill
shall be brought low: and the crooked shall be made straight, and the rough
ways plain.'' Gregory, preaching on these words, found the crooked road in the heart, bent
by injustice and made straight when it turns to the rule of justice. St
Paul in the Epistle of the day looks into the same place. The Lord who comes
``will bring to light the hidden things of darkness and will make manifest the
counsels of the hearts''; until then, ``judge not before the time''.

The Collect takes up \latin{Excita} once more, and it is the most urgent of
them:

\begin{studyframe}
\latin{Excita, quaesumus, Domine, potentiam tuam, et veni: et magna nobis
virtute succurre; ut per auxilium gratiae tuae, quod nostra peccata
praepediunt, indulgentia tuae propitiationis acceleret: Qui vivis et regnas
cum Deo Patre in unitate.}
\medskip

Bestir, O Lord, Thy might, we beseech thee, and come; and with great power come
to our aid, that, by the help of Thy grace, that which is hindered by our sins
may be hastened by Thy merciful forgiveness. Who livest.\par
\end{studyframe}

The Church admits that her sins are standing in the Lord's road, the road John
told her to make straight, and asks him to come anyway, and quickly. If you
can, let the last days before the feast be quiet.
